% ============================================================================
% CHAPTER 1 - INTRODUCTION
% Rubric mapping: CO1 (5 marks) is assessed from 1.1, 1.2, 1.3, 1.4, 1.5, 1.6.
% ============================================================================

\nomenclature{ASR}{Automatic Speech Recognition}
\nomenclature{WER}{Word Error Rate}
\nomenclature{CER}{Character Error Rate}
\nomenclature{LoRA}{Low-Rank Adaptation}
\nomenclature{LLM}{Large Language Model}
\nomenclature{VLM}{Vision-Language Model}
\nomenclature{SoM}{Set-of-Mark prompting}
\nomenclature{LOSO}{Leave-One-Speaker-Out}
\nomenclature{Banglish}{Bengali-English code-mixed speech written in the Roman alphabet}
\nomenclature{GPU}{Graphics Processing Unit}
\nomenclature{OCR}{Optical Character Recognition}

\section{Background}

Teaching in Bangladeshi universities rarely happens in a single language. In a
typical computer science class the lecturer explains an idea in Bangla, names it
in English, writes the English term on the whiteboard, and then returns to Bangla
for the worked example. A single sentence can carry both languages, as in
\textit{``erpor e ami ki korbo? not korbo''}, where the grammar is Bengali and the
technical word is English. When students write this mixture down they almost never
use the Bengali script. They type it in Roman letters, and the result is commonly
called Banglish. Students take notes in Banglish, message each other in Banglish
and search in Banglish, so Banglish is the written form that matches how the class
actually sounds.

Automatic speech recognition has become very good at the languages it was built
for. Large multilingual models such as Whisper are trained on hundreds of
thousands of hours of audio and reach usable accuracy on many languages, including
Bengali \cite{radford2023whisper}. These models, however, assume that an utterance
belongs to one language. Whisper is given a language token and a task token, and
for Bengali audio that contains English words it usually produces a fluent English
translation rather than a transcript of what was said. A Bengali speech model of
the other kind, trained only on clean Bengali, writes everything in Bengali script
and turns the English technical terms into Bengali spellings. Neither output is
the lecturer's own words, and neither can be searched by a student who writes in
Banglish.

Speech is also only half of a whiteboard lecture. In the classrooms recorded for
this work the lecturer does not use slides. Definitions, truth tables, SQL
queries, IP address calculations and diagrams appear on the board and are then
talked about. A note taken from the audio alone loses all of it. Reading the board
from the video is not a simple screenshot problem either, because the lecturer
stands in front of what he has just written, the room lights leave glare on the
surface, and the board is wiped and reused several times in one class.

Three developments make a system that handles both sides practical on ordinary
hardware. Parameter efficient fine-tuning, in particular Low-Rank Adaptation
\cite{hu2022lora}, allows a large speech model to be adapted on a few hours of new
speech by training a small number of extra weights, which fits in the memory of a
single desktop graphics card. Open vision-language models such as Qwen2.5-VL read
text inside an image and answer questions about it without any dedicated OCR stage
\cite{bai2025qwen25vl}. Instruction tuned language models such as Qwen2.5-7B can
take the transcript and the board content together and write a structured
document from them \cite{qwen2025qwen25}.

This thesis brings the three together into one system, called InsightLens. A
lecture video goes in. What comes out is a lecture note that a student can read,
containing the board redrawn without the lecturer in front of it, numbered
coloured boxes around each part of the board, an explanation written from the
lecturer's own words, and quotations of the lecturer that have been checked
against the transcript. The note is produced in two versions, one in Banglish and
one in English, so that both an English-medium reader and a student who is more
comfortable in the mixed language can use it.

\figplaceholder{fig-1-1-overview}{The InsightLens idea in one picture. A recorded
Banglish lecture is turned into a readable lecture note that contains both the
speech and the whiteboard.}{A simple left to right drawing. On the left, a camera
icon and the words ``Banglish lecture video''. In the middle, three boxes stacked
or in a row: ``Speech to Banglish text'', ``Whiteboard reconstructed and read'',
``Notes written''. On the right, a picture of the finished note page with a
coloured board image and numbered boxes. Keep it non-technical, this is the
opening figure.}{6.5cm}

\section{Rational of the Study or Motivation}

The direct motivation is that a student who misses a class, or who cannot follow
the English terms at the speed they are spoken, currently has nothing useful to
fall back on. A recording of the class exists, but three hours of video is not a
study aid. Running that video through any available transcription service returns
either English prose that was never spoken or Bengali script that no one in the
class writes. Neither helps the student revise, and neither contains the
whiteboard, which is where most of the examples live.

The second motivation is cost. Producing a correct Banglish transcript by hand is
slow work. In this project it took about one hour of human effort for every ten
minutes of video, because existing tools give the transcriber almost nothing
usable to start from. At that rate a single semester of one course is weeks of
work. If a model can be brought close enough to the lecturer's words that a human
only has to fix small errors, the cost of making lecture material accessible drops
sharply. That is a realistic goal, and it is measurable, which is why the error
rate of the speech model is treated in this thesis as a result in its own right
rather than as an internal detail.

The third motivation is that the research gap is specific and checkable. Code
switched speech recognition has been studied for Hindi-English and for several
other language pairs, and Bengali speech recognition has its own literature. What
is missing is the combination that a Bangladeshi classroom actually needs, which
is recognition of Bengali-English code-mixed speech written back in the Roman
alphabet, evaluated on real lectures rather than on read or synthetic speech.
Chapter 2 reviews the systems that come closest, including recent Bengali-English
code-switched models, and states exactly which part of this work is new and which
part is not.

Finally, the whiteboard side of the problem has a practical value beyond note
taking. A clean picture of every board a lecturer wrote, with the person removed
and the content read out as text, is useful on its own. It can be searched, it can
be pasted into a study guide, and it gives a student a record of a class that was
never written down anywhere else.

\section{Problem Statement}

The problem this thesis addresses is the following. Given an ordinary recording of
a university lecture delivered in Bengali-English code-mixed speech, in a room
with a whiteboard and no slides, produce a lecture note that is faithful to what
the lecturer said and to what the lecturer wrote.

This breaks into four smaller problems, each of which was confirmed by measurement
on our own recordings before any solution was attempted.

\textbf{The speech is not transcribed by existing systems.} On the six lectures
held out for testing in this work, off-the-shelf Whisper large-v3-turbo makes a
median character error rate of 67.7 per cent and a median word error rate of 93.9
per cent against human references. The failure is not random noise. The model
translates the speech into English and, on thirteen of the one hundred and
seventy-seven test clips, falls into a repetition loop and repeats the same phrase
until it hits the token limit. Table \ref{tab:motivating_example} shows one
segment of a digital logic class in three versions.

\begin{table}[H]
\centering
\caption{One segment of lecture BanglaASR7, as spoken and as produced by two
models. The reference is the human transcript. The off-the-shelf model translates
and then loops. The fine-tuned model writes what was said.}
\label{tab:motivating_example}
\begin{tabular}{p{2.6cm}p{10.5cm}}
\toprule
\textbf{Source} & \textbf{Text} \\
\midrule
Human reference &
\textit{done. so x-or jodi amra bujhe jai tahole last je exclusive gate sheitao
kintu amader bujha easier hoye jabe.} \\
\addlinespace
Off-the-shelf Whisper &
\textit{so, x or gate and x or gate and x or gate and x or gate and \dots}
(repeats until the token limit) \\
\addlinespace
Fine-tuned Whisper &
\textit{done. so, x or jodi amra bujhe jai, tahole last je exclusive gate, sheta o
kintu amader bujha easier hoye chabe.} \\
\bottomrule
\end{tabular}
\end{table}

\textbf{The board is hidden by the person who wrote on it.} The lecturer is never
absent from the frame. Measured over seventy-nine frames of one lecture, the
lecturer covers 23.5 per cent of the writing surface in the median frame and 7.0
per cent in the best one, and no single frame is under 5 per cent. There is
therefore no moment in the video at which a clean photograph of the board can
simply be taken.

\textbf{A summarising model writes from memory instead of from the lecture.} The
first version of this pipeline asked a language model for clear notes on the
lecture. The model produced confident textbook material, including a definition of
a NAND gate that the lecturer never gave and that was wrong. Measured against
hand-checked answer keys of what was actually written on the board, those notes
contained only 37.2 per cent of the board items across thirty-five boards. The
numbers written on the board, which a model cannot guess, were the worst category
at six items found out of sixty-seven.

\textbf{The obvious ways of scoring the system are misleading.} The measure
inherited from the earlier phase of this project, a term based F1 over a fixed
list of eighty-two English words, moves by 8.1 percentage points between two
training runs on identical data, and it moves in the wrong direction. A model that
transcribes the mixed speech correctly emits fewer English words than a model that
translates it, so the better transcript scores worse. Any evaluation built out of
English strings has this property, and it affects the note metric as well. Part of
the contribution of this thesis is to state this clearly and to measure around it.

\section{Objective}

The aim of this study is to build and evaluate an end to end system that turns a
recorded Banglish whiteboard lecture into a usable lecture note, and to report
honestly how well each part of it works. The specific objectives are listed
below.

\begin{enumerate}
  \item To build a hand-checked corpus of Bengali-English code-mixed classroom
        speech, transcribed in the Roman alphabet according to a written
        transcription standard, large enough to fine-tune and test a speech
        model.
  \item To adapt a large pretrained speech recognition model to this speech with
        Low-Rank Adaptation, and to measure the change in word and character
        error rate against the same model without adaptation, on lectures that
        were never used for training.
  \item To choose the training settings by a procedure that is written down before
        any result is seen, scored on a separate validation set, so that the
        reported test result is not the product of repeated tuning on the test
        data.
  \item To reconstruct each whiteboard from the video without the lecturer in
        front of it, without inventing any pixel, and to measure how much of the
        writing survives.
  \item To have a vision-language model read each reconstructed board, and to
        measure how much of the board's content it recovers against answer keys
        checked by hand.
  \item To generate a lecture note in both Banglish and English that is grounded
        in the board and the transcript, that points the reader at numbered
        regions of the board image, and that only quotes the lecturer using words
        that appear in the transcript.
  \item To evaluate each stage separately with paired statistical tests, to run
        ablations that show which design choices matter, and to report the results
        that did not work as clearly as the ones that did.
\end{enumerate}

\section{Methodology in Brief}

The system runs as four stages, and each stage writes a file that the next stage
reads, so any stage can be re-run or replaced on its own.

The first stage handles the speech. The audio is extracted from the video and
split into segments at the boundaries given by the transcription standard. A
pretrained Whisper large-v3-turbo model is adapted with LoRA adapters on the
training lectures, and the adapted model writes a timestamped Banglish transcript
for a lecture it has never seen.

The second stage handles the board. Frames are taken from the video at a fixed
interval. Points at which the board is wiped are detected, which splits the
lecture into eras, one board per era. Within an era the frame is cut into
overlapping tiles, and each tile is taken from a moment when the lecturer was not
standing in front of that particular tile. The lecturer is located by a pretrained
person segmentation network together with a shadow mask. The tiles are blended
back into one image, so every pixel of the result is unmodified camera output.

The third stage reads the board. Our own code finds the blocks of writing on the
clean board, draws a numbered coloured box around each one, and a vision-language
model is asked to name and transcribe each numbered box. This is the Set-of-Mark
style of prompting \cite{yang2023som}, where the marks are drawn on the image and
the model answers per mark.

The fourth stage writes the note. For each board, a language model is given that
board's boxes, the board text, and the part of the transcript that belongs to that
board, and is asked for one section of notes that refers to the boxes by number
and colour. A checker then deletes any quotation that does not appear word for
word in the transcript, and any reference to a box number that does not exist. The
sections are assembled into a page in Markdown and in HTML with the board image
embedded.

Evaluation is done per stage and never by eye alone. Speech is scored by per-clip
median word and character error rate against human transcripts, with paired
Wilcoxon signed-rank tests and sign tests. The board reconstruction is scored by
the fraction of tiles that found a clear view and by a human check of whether any
writing was lost. Board reading and note quality are scored by recall of items
from hand-verified answer keys of what each board contained. Chapter 4 gives the
full detail of each stage and Chapter 5 gives the results.

\section{Scopes and Challenges}

\textbf{Scope.} The dataset is forty-four lecture recordings from three lecturers
at one university, 8.33 hours of video in total, of which twenty-eight lectures
covering 5.15 hours of speech carry hand-checked transcripts. The subjects are
introductory programming, digital logic, databases and computer networking. Every
recording is of a whiteboard class with no slides, taken with a fixed camera. The
notes are produced in Banglish and in English. Notes in the Bengali script were
considered and dropped, because they would be produced by translation alone and
would add nothing that the system itself contributes. The system is built for
recorded lectures, not for live captioning.

\textbf{What this thesis does not claim.} It does not claim that the speech model
is accurate enough to be used without a human check, and Chapter 5 states the
remaining error rate plainly. It does not claim that the generated notes are good
to read, because that has not been tested with readers. The measure used for the
notes counts whether board content reached the page, which is a real and checkable
property, but it says nothing about whether the surrounding explanation is well
written or even correct. A reader study is described in Chapter 6 as the next
step.

\textbf{Challenges.} Five difficulties shaped the work.

\begin{enumerate}
  \item \textbf{No standard spelling.} Banglish has no agreed orthography, so the
        same word can be written \textit{ami}, \textit{aami} or \textit{amee}. A
        written transcription guide was produced and enforced by a validator, and
        a spelling-tolerant version of the error rate was computed to show how
        much of the error is spelling and how much is genuine.
  \item \textbf{No dataset existed.} There is no public corpus of Bengali-English
        code-mixed classroom speech in Roman script. The corpus had to be
        recorded and transcribed by hand, and this was the single largest cost in
        the project.
  \item \textbf{Limited hardware.} All training was done on one desktop with a
        12 GB graphics card and all vision-language work on a second desktop with
        a 32 GB card. This ruled out the largest vision-language models, and the
        thesis says so rather than pretending the comparison was a free choice.
  \item \textbf{Unstable training.} The configuration chosen by the tuning
        procedure diverged on one of its two random seeds and produced a model
        worse than doing nothing. This is reported in full in Chapter 5, together
        with the settings that were stable.
  \item \textbf{Metrics that mislead.} As described in the problem statement,
        measures built out of English strings reward a model that translates. Two
        separate results in this thesis are affected by it, and both are reported
        with the explanation beside them.
\end{enumerate}

\section{Team Overview}

This thesis was carried out by a team of three students of the Department of
Computer Science and Engineering, Brac University, under the supervision of
Dr. Md. Golam Rabiul Alam and the co-supervision of Md. Tanzim Reza.

% TODO (Rafi): confirm or correct the division of work in the paragraph below
% before submission. Student IDs go on the title page and the declaration page.
The work was divided into three areas. Data collection and transcription was the
largest shared task, since every transcript in the corpus was checked word by word
against the audio by a member of the team. Model training and evaluation covered
the speech recognition experiments, the tuning procedure and the statistical
testing. The vision and note generation side covered board reconstruction, the
vision-language prompting and the note builder. All three members took part in the
verification of the board answer keys, which had to be done by a person looking at
the board image without seeing any model output, and in the writing of this
report.
